\documentclass{article}
\usepackage[utf8]{inputenc}
\usepackage{amsmath,amssymb}
\usepackage[most]{tcolorbox}
\usepackage{geometry}
\geometry{margin=2.5cm}

\newcommand{\step}[1]{\par\noindent\hangindent=3.5em\hangafter=1%
  \makebox[3.5em][l]{\textbf{#1.}}}
\newcommand{\steptag}[1]{\hfill{\small\textsf{[#1]}}}

\newtcolorbox{proofbox}[1]{
  colback=gray!5, colframe=gray!60,
  fonttitle=\bfseries, title={#1},
  breakable, enhanced,
  left=4pt, right=4pt, top=4pt, bottom=4pt,
  before skip=10pt, after skip=10pt
}

\begin{document}

\begin{proofbox}{Corrected th\_almost\_idemp audit: let H be nonzero and let...}
\step{1} Corrected th\_almost\_idemp audit: let H be nonzero and let Phi:B(H)->B(H) be UCP with ||Phi\textasciicircum{}2-Phi||\_cb <= eta < 1/4; for tilde-Phi=(1/2)(I+(2Phi-I)(I-4(Phi-Phi\textasciicircum{}2))\textasciicircum{}(-1/2)), A=Im(tilde-Phi), and X star Y=tilde-Phi(XY), one has tilde-Phi\textasciicircum{}2=tilde-Phi, both amplified associativity identities Phi\_assoc1 and Phi\_assoc2 have error at most 10*eta*||X||||Y||||Z|| after the local source type/index corrections, and for sufficiently small universal eta the inherited operator-space norms, involution, and unit make A an extended epsilon\_AI(eta)-C*-algebra with epsilon\_AI=max\{r,20eta+2(M\textasciicircum{}5-1),3r-r\textasciicircum{}2\}=O(eta), r=(3/2)((1-4eta)\textasciicircum{}(-1/2)-1), M=1+r, dimension-free at every amplification. \steptag{validated} \\[6pt]
\step{1.1} Spectral-correction package: in the Banach algebra of completely bounded maps under composition, F=Phi-Phi\textasciicircum{}2 has norm\_cb at most eta, A0=2Phi-I satisfies A0\textasciicircum{}2=I-4F, B=(I-4F)\textasciicircum{}(-1/2) is given by the convergent binomial series, S=A0 B satisfies S\textasciicircum{}2=I, and tilde-Phi=(I+S)/2 is an idempotent unital star-preserving map with ||tilde-Phi-Phi||\_cb at most r=(3/2)((1-4eta)\textasciicircum{}(-1/2)-1); consequently its image is a closed complete self-adjoint operator space containing I, and r=O(eta) universally. \steptag{validated} \\[6pt]
\step{1.1.1} Functional calculus calculation: set F=Phi-Phi\textasciicircum{}2 and A0=2Phi-I. Since F is a polynomial in Phi, A0 and F commute, A0\textasciicircum{}2=I-4F, and ||4F||\_cb<1. The norm-convergent binomial series B=sum\_\{k>=0\} binom(2k,k)F\textasciicircum{}k satisfies B=(I-4F)\textasciicircum{}(-1/2), B\textasciicircum{}2=(I-4F)\textasciicircum{}(-1), and commutes with A0. Thus S=A0B has S\textasciicircum{}2=I and P=(I+S)/2 has P\textasciicircum{}2=P. Moreover P-Phi=(1/2)A0(B-I), ||A0||\_cb<=2||Phi||\_cb+1<=3, and positivity of the scalar series coefficients gives ||B-I||\_cb<=sum\_\{k>=1\} binom(2k,k)eta\textasciicircum{}k=(1-4eta)\textasciicircum{}(-1/2)-1, hence ||P-Phi||\_cb<=r. \steptag{validated} \\[4pt]
\step{1.1.2} Structure and asymptotics: UCP implies Phi(I)=I and Phi(X\textasciicircum{}*)=Phi(X)\textasciicircum{}*. Therefore F(I)=0, B(I)=I, A0(I)=I and P(I)=I; the real-coefficient series and composition show P(X\textasciicircum{}*)=P(X)\textasciicircum{}*. A bounded idempotent has Im(P)=Ker(I-P), so its image is closed, complete in the inherited norm, self-adjoint, contains I, and its restricted matrix norms form the inherited complete operator-space structure. The same formulas commute with every matrix amplification. Finally, for 0<=eta<=1/8 the mean-value theorem gives r<=6 sqrt(2) eta, so r=O(eta) with a universal dimension-free constant. \steptag{validated} \\[4pt]
\step{1.2} First amplified associativity identity: for every n and T=I\_Mn tensor Phi on B(C\textasciicircum{}n tensor H), and all X,Y,Z at that level, ||T(T(T(X)T(Y))T(Z))-T(T(X)T(Y)T(Z))|| is at most 10 eta ||X||||Y||||Z||; the proof uses an internal, explicitly typed Stinespring-stack rectangle factorization and only UCP plus ||T\textasciicircum{}2-T|| at most eta. \steptag{validated} \\[6pt]
\step{1.2.1} Elementary uniform defect ledger: at a fixed matrix level T is UCP, hence ||T||<=1 and T(U\textasciicircum{}*)=T(U)\textasciicircum{}*, while D=T\textasciicircum{}2-T has ||D||<=eta. Thus ||T\textasciicircum{}2(U)-T(U)||<=eta||U|| and, by inserting T(U)T\textasciicircum{}2(V), ||T\textasciicircum{}2(U)T\textasciicircum{}2(V)-T(U)T(V)||<=2eta||U||||V||. Consequently G\_X=T\textasciicircum{}2(T(X\textasciicircum{}*)T(X))-T(T\textasciicircum{}2(X\textasciicircum{}*)T\textasciicircum{}2(X)) has ||G\_X||<=3eta||X||\textasciicircum{}2. These estimates use no dimension, trace, Kraus sum, or ancillary basis. \steptag{validated} \\[4pt]
\step{1.2.2} Internal Stinespring rectangle and correction ledger: the iterative GNS/Stinespring construction for UCP T gives Hilbert spaces K\_j, isometries V\_j:K\_\{j-1\}->K\_j, and contractive unital star-homomorphisms u\_j:B(K\_\{j-1\})->B(K\_j), with the one-step compression identity u\_j(V\_j\textasciicircum{}*UV\_j)=V\_\{j+1\}\textasciicircum{}*u\_\{j+1\}(U)V\_\{j+1\} for U in B(K\_j), and its iterates. All four source memberships are therefore U in B(K\_j), never U in K\_j. For R\_k(X)=u\_\{2+k,1\}((I\_1-V\_1V\_1\textasciicircum{}*)u\_1(T(X)))V\_\{2+k\}V\_\{1+k\}, exact multiplication and the iterated compression identity give R\_k(X)\textasciicircum{}*R\_k(X)=u\_\{k,0\}(G\_X), hence ||R\_k(X)||<=sqrt(3eta)||X||. The auxiliary identity has final factor V\_\{1+k\}: V\_\{1+k\}\textasciicircum{}*V\_\{2+k\}\textasciicircum{}*u\_\{2+k,1\}(Z)V\_\{2+k\}V\_\{1+k\}=V\_\{1+k\}\textasciicircum{}*u\_\{1+k,0\}(V\_1\textasciicircum{}*ZV\_1)V\_\{1+k\}; using V\_1 there is ill-typed. With A=T(X), B=T(Y), C=T(Z), A1=T\textasciicircum{}2(X), B1=T\textasciicircum{}2(Y), C1=T\textasciicircum{}2(Z), the operator W=V\_1\textasciicircum{}*R\_1(X\textasciicircum{}*)\textasciicircum{}*u\_\{3,0\}(T(Y))V\_3R\_0(Z) obeys ||W||<=3eta||X||||Y||||Z||, and expanding both I\_1-V\_1V\_1\textasciicircum{}* factors then compressing gives exactly W=T(Q), Q=T(T(AB)C)-T\textasciicircum{}2(AB)C1-T(A1B1C)+T(A1B1)C1. \steptag{validated} \\[6pt]
\step{1.2.2.1} Explicit recursive dilation, proved from complete positivity rather than imported: for any Hilbert space K and UCP map Psi:B(K)->B(K), put on the algebraic tensor product B(K) tensor K the semidefinite form <sum\_i A\_i tensor xi\_i, sum\_j B\_j tensor eta\_j>=sum\_ij <xi\_i,Psi(A\_i\textasciicircum{}*B\_j)eta\_j>. Complete positivity of Psi makes this form positive because [A\_i\textasciicircum{}*A\_j] is a positive operator matrix. After quotienting its null space and completing, define pi(C)[A tensor xi]=[CA tensor xi] and Vxi=[I tensor xi]. The inequality C\textasciicircum{}*C<=||C||\textasciicircum{}2 I makes pi(C) bounded with norm at most ||C||; the defining form gives pi(C)\textasciicircum{}*=pi(C\textasciicircum{}*), pi(CD)=pi(C)pi(D), and pi(I)=I. Unitality gives V\textasciicircum{}*V=I, and direct evaluation gives V\textasciicircum{}*pi(C)V=Psi(C). Apply this construction first to Psi\_0=T on B(K\_0), obtaining (K\_1,u\_1,V\_1), and recursively to the UCP map Psi\_j:B(K\_j)->B(K\_j), Psi\_j(U)=u\_j(V\_j\textasciicircum{}*UV\_j), obtaining (K\_\{j+1\},u\_\{j+1\},V\_\{j+1\}). Thus each V\_j:K\_\{j-1\}->K\_j is an isometry, each u\_j:B(K\_\{j-1\})->B(K\_j) is a contractive unital star-homomorphism, T(H)=V\_1\textasciicircum{}*u\_1(H)V\_1, and for every j>=1 and U in B(K\_j), u\_j(V\_j\textasciicircum{}*UV\_j)=V\_\{j+1\}\textasciicircum{}*u\_\{j+1\}(U)V\_\{j+1\}. This proves the claimed iterative system with the source type U in B(K\_j), without an external Stinespring axiom. \steptag{validated} \\[4pt]
\step{1.2.2.2} Self-contained rectangle and defect calculation: let T:B(K\_0)->B(K\_0) be UCP with ||T\textasciicircum{}2-T||<=eta and let (K\_j,u\_j,V\_j) be any iterative system furnished by the preceding explicit construction. For b>a define u\_\{b,a\}=u\_b o ... o u\_\{a+1\} and u\_\{a,a\}=id; write I\_j for the identity of K\_j and P\_1=V\_1V\_1\textasciicircum{}*. Induction in the one-step compression identity gives, for k>=0, H in B(K\_0), and Z in B(K\_1), (i) V\_\{k+1\}\textasciicircum{}*u\_\{k+1,0\}(H)V\_\{k+1\}=u\_\{k,0\}(T(H)), and (ii) V\_\{k+2\}\textasciicircum{}*u\_\{k+2,1\}(Z)V\_\{k+2\}=u\_\{k+1,0\}(V\_1\textasciicircum{}*ZV\_1). Hence the fully typed auxiliary identity is V\_\{k+1\}\textasciicircum{}*V\_\{k+2\}\textasciicircum{}*u\_\{k+2,1\}(Z)V\_\{k+2\}V\_\{k+1\}=V\_\{k+1\}\textasciicircum{}*u\_\{k+1,0\}(V\_1\textasciicircum{}*ZV\_1)V\_\{k+1\}; its final factor must be V\_\{k+1\}, since V\_1 has domain K\_0 rather than K\_k when k>0. Put E\_X=(I\_1-P\_1)u\_1(T(X)) in B(K\_1) and H\_X=V\_1\textasciicircum{}*E\_X\textasciicircum{}*E\_XV\_1. Multiplicativity of u\_1 and T(U)=V\_1\textasciicircum{}*u\_1(U)V\_1 give exactly H\_X=T(T(X\textasciicircum{}*)T(X))-T\textasciicircum{}2(X\textasciicircum{}*)T\textasciicircum{}2(X). Define G\_X=T(H\_X)=T\textasciicircum{}2(T(X\textasciicircum{}*)T(X))-T(T\textasciicircum{}2(X\textasciicircum{}*)T\textasciicircum{}2(X)). Contractivity and the defect estimate give ||T\textasciicircum{}2(T(X\textasciicircum{}*)T(X))-T(T(X\textasciicircum{}*)T(X))||<=eta||X||\textasciicircum{}2 and ||T(T(X\textasciicircum{}*)T(X)-T\textasciicircum{}2(X\textasciicircum{}*)T\textasciicircum{}2(X))||<=2eta||X||\textasciicircum{}2, the latter by adding and subtracting T(X\textasciicircum{}*)T\textasciicircum{}2(X); therefore ||G\_X||<=3eta||X||\textasciicircum{}2. For R\_k(X)=u\_\{k+2,1\}(E\_X)V\_\{k+2\}V\_\{k+1\}, exact multiplication followed by (ii) and then (i) yields R\_k(X)\textasciicircum{}*R\_k(X)=V\_\{k+1\}\textasciicircum{}*u\_\{k+1,0\}(H\_X)V\_\{k+1\}=u\_\{k,0\}(T(H\_X))=u\_\{k,0\}(G\_X). Since star-homomorphisms are contractive, ||R\_k(X)||\textasciicircum{}2<=||G\_X||, so ||R\_k(X)||<=sqrt(3eta)||X||. \steptag{validated} \\[4pt]
\step{1.2.2.3} Exact W expansion with all symbols declared: under the same iterative system, set A=T(X), B=T(Y), C=T(Z), A\_1=T(A)=T\textasciicircum{}2(X), B\_1=T(B)=T\textasciicircum{}2(Y), C\_1=T(C)=T\textasciicircum{}2(Z), P\_1=V\_1V\_1\textasciicircum{}*, and E\_U=(I\_1-P\_1)u\_1(T(U)). Then R\_0(Z)=u\_\{2,1\}(E\_Z)V\_2V\_1 and R\_1(X\textasciicircum{}*)=u\_\{3,1\}(E\_\{X\textasciicircum{}*\})V\_3V\_2, with E\_\{X\textasciicircum{}*\}\textasciicircum{}*=u\_1(A)(I\_1-P\_1). Put D=V\_1\textasciicircum{}*u\_1(A)(I\_1-P\_1)u\_1(B)V\_1=T(AB)-A\_1B\_1. For W=V\_1\textasciicircum{}*R\_1(X\textasciicircum{}*)\textasciicircum{}*u\_\{3,0\}(B)V\_3R\_0(Z), multiplicativity and the two compression identities give the explicit chain W=V\_1\textasciicircum{}*V\_2\textasciicircum{}*u\_2(u\_1(D)(I\_1-P\_1)u\_1(C))V\_2V\_1=T(V\_1\textasciicircum{}*u\_1(D)(I\_1-P\_1)u\_1(C)V\_1)=T(T(DC)-T(D)C\_1). Substituting D and using linearity yields W=T(Q), where Q=T(T(AB)C)-T\textasciicircum{}2(AB)C\_1-T(A\_1B\_1C)+T(A\_1B\_1)C\_1, exactly as claimed. Finally, the rectangle bound from the preceding calculation and contractivity of u\_\{3,0\} give ||W||<=||R\_1(X\textasciicircum{}*)|| ||B|| ||R\_0(Z)||<=sqrt(3eta)||X|| ||Y|| sqrt(3eta)||Z||=3eta||X||||Y||||Z||. Thus neither W=T(Q) nor its norm estimate is an undeclared premise. \steptag{validated} \\[6pt]
\step{1.2.2.3.1} Local derivation of the two rectangle estimates, with no sibling imported. Let U be any operator on K\_0. From the one-step compression identity of the iterative system, induction gives V\_\{k+2\}\textasciicircum{}*u\_\{k+2,1\}(Z)V\_\{k+2\}=u\_\{k+1,0\}(V\_1\textasciicircum{}*ZV\_1) for Z in B(K\_1), and V\_\{k+1\}\textasciicircum{}*u\_\{k+1,0\}(H)V\_\{k+1\}=u\_\{k,0\}(T(H)) for H in B(K\_0). Put E\_U=(I\_1-P\_1)u\_1(T(U)), H\_U=V\_1\textasciicircum{}*E\_U\textasciicircum{}*E\_UV\_1, and G\_U=T(H\_U). Multiplicativity and T(H)=V\_1\textasciicircum{}*u\_1(H)V\_1 yield H\_U=T(T(U\textasciicircum{}*)T(U))-T\textasciicircum{}2(U\textasciicircum{}*)T\textasciicircum{}2(U), hence G\_U=T\textasciicircum{}2(T(U\textasciicircum{}*)T(U))-T(T\textasciicircum{}2(U\textasciicircum{}*)T\textasciicircum{}2(U)). Writing V=T(U\textasciicircum{}*)T(U), the first difference [T\textasciicircum{}2(V)-T(V)] has norm at most eta||U||\textasciicircum{}2. Moreover T(U\textasciicircum{}*)T(U)-T\textasciicircum{}2(U\textasciicircum{}*)T\textasciicircum{}2(U)=T(U\textasciicircum{}*)[T(U)-T\textasciicircum{}2(U)]+[T(U\textasciicircum{}*)-T\textasciicircum{}2(U\textasciicircum{}*)]T\textasciicircum{}2(U), so contractivity of T and T\textasciicircum{}2 and ||T\textasciicircum{}2-T||<=eta bound the norm of its image under T by 2eta||U||\textasciicircum{}2. Thus ||G\_U||<=3eta||U||\textasciicircum{}2. For R\_k(U)=u\_\{k+2,1\}(E\_U)V\_\{k+2\}V\_\{k+1\}, exact multiplication followed by the two displayed compression identities gives R\_k(U)\textasciicircum{}*R\_k(U)=u\_\{k,0\}(G\_U). Contractivity of the star-homomorphism u\_\{k,0\} therefore gives ||R\_k(U)||<=sqrt(3eta)||U||. In particular ||R\_1(X\textasciicircum{}*)||<=sqrt(3eta)||X|| and ||R\_0(Z)||<=sqrt(3eta)||Z||. \steptag{validated} \\[4pt]
\step{1.2.2.3.2} Complete W conclusion using the locally proved rectangle estimates. With the target node notation, multiplicativity and the two compression identities give D=V\_1\textasciicircum{}*u\_1(A)(I\_1-P\_1)u\_1(B)V\_1=T(AB)-A\_1B\_1 and W=V\_1\textasciicircum{}*V\_2\textasciicircum{}*u\_2(u\_1(D)(I\_1-P\_1)u\_1(C))V\_2V\_1=T(V\_1\textasciicircum{}*u\_1(D)(I\_1-P\_1)u\_1(C)V\_1)=T(T(DC)-T(D)C\_1). Substitution of D and linearity give W=T(Q), Q=T(T(AB)C)-T\textasciicircum{}2(AB)C\_1-T(A\_1B\_1C)+T(A\_1B\_1)C\_1. By the preceding child, contractivity of T and u\_\{3,0\}, ||B||<=||Y||, and the isometry norms, ||W||<=||R\_1(X\textasciicircum{}*)|| ||B|| ||R\_0(Z)||<=3eta||X||||Y||||Z||. Thus both the exact expansion and the norm conclusion follow without importing the sibling rectangle node. \steptag{validated} \\[4pt]
\step{1.2.3} Five-plus-two eta reduction and conclusion: for F1=T(T(AB)C)-T(ABC), the last two paired terms in Q satisfy ||T(A1B1)-T\textasciicircum{}2(AB)||<=||A1B1-AB||+||T(AB)-T\textasciicircum{}2(AB)||<=3eta||X||||Y||, while ||T(A1B1C)-T(ABC)||<=2eta||X||||Y||||Z||. Hence ||Q-F1||<=5eta||X||||Y||||Z||. Contractivity gives ||T(Q)-T(F1)||<=5eta product, and because F1 is a difference of two T-valued terms, the defect estimate applied once to each gives ||T(F1)-F1||<=2eta product. Therefore ||W-F1||<=7eta product; together with ||W||<=3eta product this yields ||F1||<=10eta||X||||Y||||Z||, exactly Phi\_assoc1. \steptag{validated} \\[6pt]
\step{1.2.3.1} Fixed-level scope and rectangle input: fix n and X,Y,Z in B(C\textasciicircum{}n tensor H), set T=I\_Mn tensor Phi, A=T(X), B=T(Y), C=T(Z), A1=T\textasciicircum{}2(X), B1=T\textasciicircum{}2(Y), C1=T\textasciicircum{}2(Z), and p=||X||||Y||||Z||. By node 1.2.1, ||T||<=1 and ||T\textasciicircum{}2(U)-T(U)||<=eta||U||. Using the Stinespring objects V\_j,u\_j,R\_k defined in node 1.2.2, set W=V\_1\textasciicircum{}* R\_1(X\textasciicircum{}*)\textasciicircum{}* u\_\{3,0\}(T(Y)) V\_3 R\_0(Z) and Q=T(T(AB)C)-T\textasciicircum{}2(AB)C1-T(A1B1C)+T(A1B1)C1. The exact expansion in node 1.2.2 gives W=T(Q) and ||W||<=3eta p. \steptag{validated} \\[6pt]
\step{1.2.3.1.1} Self-contained exact expansion, independent of pending node 1.2.2: in the fixed-level notation of 1.2.3.1 use the validated recursive dilation of 1.2.2.1, write u\_\{b,a\}=u\_b o ... o u\_\{a+1\}, P\_1=V\_1V\_1\textasciicircum{}*, E\_U=(I\_1-P\_1)u\_1(T(U)), R\_0(Z)=u\_\{2,1\}(E\_Z)V\_2V\_1, and R\_1(X\textasciicircum{}*)=u\_\{3,1\}(E\_\{X\textasciicircum{}*\})V\_3V\_2. Star preservation gives E\_\{X\textasciicircum{}*\}\textasciicircum{}*=u\_1(A)(I\_1-P\_1). Set D=V\_1\textasciicircum{}*u\_1(A)(I\_1-P\_1)u\_1(B)V\_1=T(AB)-A\_1B\_1. Multiplicativity and the validated compression identity first give W=V\_1\textasciicircum{}*V\_2\textasciicircum{}*u\_2(u\_1(D)(I\_1-P\_1)u\_1(C))V\_2V\_1, then W=T(V\_1\textasciicircum{}*u\_1(D)(I\_1-P\_1)u\_1(C)V\_1)=T(T(DC)-T(D)C\_1). Since D=T(AB)-A\_1B\_1, linearity gives W=T(T(T(AB)C)-T\textasciicircum{}2(AB)C\_1-T(A\_1B\_1C)+T(A\_1B\_1)C\_1)=T(Q). Every product is typed in the operator algebras B(K\_j). \steptag{validated} \\[4pt]
\step{1.2.3.1.2} Self-contained norm estimate using only validated nodes: node 1.2.2.2 proves for the same explicitly defined rectangles that ||R\_k(U)||<=sqrt(3eta)||U||. Since V\_1 and V\_3 are isometries and u\_\{3,0\} is contractive, ||W||<=||R\_1(X\textasciicircum{}*)|| ||u\_\{3,0\}(B)|| ||R\_0(Z)||<=sqrt(3eta)||X|| ||B|| sqrt(3eta)||Z||. Node 1.2.1 gives ||B||=||T(Y)||<=||Y||; hence ||W||<=3eta||X||||Y||||Z||=3eta p. This does not invoke the pending parent 1.2.2 or pending sibling 1.2.2.3. \steptag{validated} \\[4pt]
\step{1.2.3.2} Elementary perturbation estimates: contractivity of T and T\textasciicircum{}2 gives ||A||,||A1||<=||X||, ||B||,||B1||<=||Y||, and ||C||,||C1||<=||Z||. Also ||A1-A||<=eta||X|| and ||B1-B||<=eta||Y||. Hence A1B1-AB=(A1-A)B1+A(B1-B), so ||A1B1-AB||<=2eta||X||||Y||. Consequently ||T(A1B1C)-T(ABC)||<=||(A1B1-AB)C||<=2eta p, and ||T(A1B1)-T\textasciicircum{}2(AB)||<=||T(A1B1)-T(AB)||+||T(AB)-T\textasciicircum{}2(AB)||<=2eta||X||||Y||+eta||AB||<=3eta||X||||Y||. \steptag{validated} \\[4pt]
\step{1.2.3.3} Exact five-eta decomposition: with F1=T(T(AB)C)-T(ABC), direct cancellation in the displayed definition of Q gives Q-F1=[T(ABC)-T(A1B1C)]+[T(A1B1)-T\textasciicircum{}2(AB)]C1. Therefore the preceding estimates and ||C1||<=||Z|| give ||Q-F1||<=2eta p+3eta p=5eta p. \steptag{validated} \\[4pt]
\step{1.2.3.4} Two-eta defect correction and conclusion: write R=T(AB)C and S=ABC, so F1=T(R)-T(S). Then T(F1)-F1=[T\textasciicircum{}2(R)-T(R)]-[T\textasciicircum{}2(S)-T(S)], whence ||T(F1)-F1||<=eta(||R||+||S||)<=2eta p because ||T(AB)||<=||A||||B|| and ||A||<=||X||, ||B||<=||Y||, ||C||<=||Z||. Since W=T(Q), contractivity and the five-eta bound give ||W-F1||<=||T(Q)-T(F1)||+||T(F1)-F1||<=5eta p+2eta p=7eta p. Finally ||F1||<=||W||+||W-F1||<=10eta p; substituting A=T(X), B=T(Y), C=T(Z) identifies F1 exactly with the Phi\_assoc1 difference in node 1.2. \steptag{validated} \\[4pt]
\step{1.3} Adjoint reduction for the second identity: for every star-preserving T, the first associativity estimate applied to (Z\textasciicircum{}*,Y\textasciicircum{}*,X\textasciicircum{}*) and then adjointed is exactly ||T(T(X)T(T(Y)T(Z)))-T(T(X)T(Y)T(Z))|| at most the same constant times ||X||||Y||||Z||; hence Phi\_assoc2 follows from Phi\_assoc1 with constant 10 at every amplification. \steptag{validated} \\[4pt]
\step{1.4} Extended approximate C-star package: writing P=tilde-Phi, d=||P-Phi||\_cb at most r and M=1+r, the inherited matrix norms, involution and unit and the product X star Y=P(XY) make every M\_n tensor Im(P) an epsilon\_AI-algebra with exact unit and involution, submultiplicativity error r, associativity error at most 20eta+2(M\textasciicircum{}5-1), and C-star lower-norm error at most 3r-r\textasciicircum{}2; for sufficiently small universal eta their maximum epsilon\_AI is less than 1, is O(eta), and all constants are independent of n and dim(H). \steptag{validated} \\[6pt]
\step{1.4.1} All nonassociative axioms with explicit errors: at each matrix level let T=I\_Mn tensor Phi and P=I\_Mn tensor tilde-Phi, so ||P-T||<=r and ||P||<=M. The range of P is a closed complete self-adjoint operator space; P(I)=I, P is star-preserving, and P\textasciicircum{}2=P. Hence X star Y=P(XY) is bilinear and range-valued, I is an exact unit, and (X star Y)\textasciicircum{}*=Y\textasciicircum{}* star X\textasciicircum{}*. Also ||X star Y||<=M||X||||Y||, giving submultiplicativity error r. For X in Im(P), ||T(X)||>=||X||-r||X||=(1-r)||X||; UCP Schwarz gives ||T(X\textasciicircum{}*X)||>=||T(X)||\textasciicircum{}2, and therefore ||P(X\textasciicircum{}*X)||>=||T(X\textasciicircum{}*X)||-r||X|$|\textasciicircum{}2\rangle$=(1-3r+r\textasciicircum{}2)||X||\textasciicircum{}2. Thus the C-star lower-norm error is 3r-r\textasciicircum{}2 for small r. \steptag{validated} \\[4pt]
\step{1.4.2} Associativity, amplification, and packaging under def-extended-epsilon-cstar-algebra: put L\_P=P(P(P(X)P(Y))P(Z)) and R\_P=P(P(X)P(P(Y)P(Z))). Telescoping replacement of the five map occurrences, using ||T||<=1, ||P||<=M and ||P-T||<=r=M-1, gives ||L\_P-L\_T|| and ||R\_P-R\_T|| each at most (M\textasciicircum{}5-1)||X||||Y||||Z||. Phi\_assoc1 and Phi\_assoc2 put L\_T and R\_T within 10eta product of the same T(T(X)T(Y)T(Z)), so ||L\_P-R\_P||<=(20eta+2(M\textasciicircum{}5-1)) product. For X,Y,Z in Im(P), L\_P=(X star Y) star Z and R\_P=X star (Y star Z). Taking the maximum of this error, r, and 3r-r\textasciicircum{}2 verifies every epsilon-C-star axiom at every matrix level, and the complete self-adjoint operator-space, multiplication and unit therefore satisfy def-extended-epsilon-cstar-algebra. Since r=O(eta), M\textasciicircum{}5-1<=31r for 0<=r<=1, and all displayed errors tend to zero, a universal eta\_0>0 makes epsilon\_AI<1 and epsilon\_AI=O(eta), uniformly in n and dim(H). \steptag{validated} \\[6pt]
\step{1.4.2.1} Dependency-explicit bridge. Require validated nodes 1.2, 1.3, and 1.4.1. Fix an amplification n and abbreviate q=||X||||Y||||Z||. Node 1.2 gives ||L\_T-C\_T||<=10 eta q for L\_T=T(T(T(X)T(Y))T(Z)) and C\_T=T(T(X)T(Y)T(Z)); node 1.3 gives ||R\_T-C\_T||<=10 eta q for R\_T=T(T(X)T(T(Y)T(Z))). For either of the five-occurrence expressions L or R, replace its occurrences of P one at a time by T. Multilinearity of composition and multiplication, ||T||<=1, ||P||<=M, and ||P-T||<=r show that the term in which the j-th occurrence is replaced has norm at most r M\textasciicircum{}j q for some j in \{0,1,2,3,4\}; hence each telescoping sum is at most r(1+M+M\textasciicircum{}2+M\textasciicircum{}3+M\textasciicircum{}4)q=(M\textasciicircum{}5-1)q because M=1+r. Therefore ||L\_P-R\_P||<=||L\_P-L\_T||+||L\_T-C\_T||+||C\_T-R\_T||+||R\_T-R\_P||<=(20 eta+2(M\textasciicircum{}5-1))q. If X,Y,Z lie in Im(P), idempotence of P identifies L\_P=(X star Y) star Z and R\_P=X star (Y star Z). Node 1.4.1 supplies at this same arbitrary matrix level the complete self-adjoint operator-space structure, range-valued bilinear multiplication, exact unit and involution, submultiplicativity error r, and C-star lower-norm error 3r-r\textasciicircum{}2. Thus, with epsilon\_AI=max\{r,20 eta+2(M\textasciicircum{}5-1),3r-r\textasciicircum{}2\}, every amplification has all the required epsilon-C-star properties, so the common operator space, multiplication, and unit satisfy def-extended-epsilon-cstar-algebra. Finally r=(3/2)((1-4 eta)\textasciicircum{}(-1/2)-1)=O(eta) as eta tends to 0; when 0<=r<=1, M\textasciicircum{}5-1=(1+r)\textasciicircum{}5-1<=31r. Hence epsilon\_AI=O(eta) and tends to 0, so a universal eta\_0>0 makes epsilon\_AI<1. All input estimates are amplification-uniform, so the conclusion is independent of n and dim(H). \steptag{validated} \\[4pt]
\end{proofbox}

\end{document}
